\documentclass[aspectratio=169]{beamer}

\usepackage{listings}
\usepackage{xcolor}

\definecolor{codegray}{rgb}{0.5,0.5,0.5}
\definecolor{codepurple}{rgb}{0.58,0,0.82}
\definecolor{backcolour}{rgb}{0.95,0.95,0.92}

\lstdefinestyle{customc}{
  backgroundcolor=\color{backcolour},
  commentstyle=\color{codegray}\itshape,
  keywordstyle=\color{blue}\bfseries,
  numberstyle=\tiny\color{gray},
  stringstyle=\color{codepurple},
  basicstyle=\ttfamily\footnotesize,
  breakatwhitespace=false,
  breaklines=true,
  captionpos=b,
  keepspaces=true,
  numbers=left,
  numbersep=5pt,
  showspaces=false,
  showstringspaces=false,
  showtabs=false,
  tabsize=2,
  frame=single,
  language=C++
}

\lstset{style=customc}

\usepackage[spanish]{babel}
\usepackage[utf8]{inputenc}
\usepackage[T1]{fontenc}
\usepackage{lmodern}
\usepackage{amsmath, amsfonts}
\usepackage{hyperref}

\usetheme{Madrid}

\title{Búsqueda exhaustiva y backtracking}
\author{Sebastian Mestre}
\date{}

\begin{document}

\begin{frame}
  \titlepage
\end{frame}

\section{Fuerza bruta}

\begin{frame}{¿Cuándo alcanza con fuerza bruta?}
  \begin{itemize}
    \item Para problemas con entrada pequeña (por ejemplo \(N < 20\)) podemos probar \textbf{todas las posibilidades}.
    \item La idea es explorar el \textbf{espacio de soluciones completo} y quedarnos con la mejor / alguna válida.
    \item Aunque parezca simple, muchos equipos se complican al implementarlo y pierden puntos en problemas fáciles.
    \item Objetivo de la clase:
      \begin{itemize}
        \item Ver formas estándar de implementar fuerza bruta eficientemente.
        \item Aprender trucos para \textbf{escribir poco código} y evitar errores.
      \end{itemize}
  \end{itemize}
\end{frame}

\section{Subconjuntos y bitmasks}

\begin{frame}{Fuerza bruta sobre subconjuntos}
  \begin{itemize}
    \item Para un conjunto de \(n\) elementos, hay \(2^n\) subconjuntos posibles.
    \item Podemos representar cada subconjunto con un entero de \(0\) a \(2^n - 1\):
      \begin{itemize}
        \item El bit \(i\) indica si el elemento \(i\) está presente en el subconjunto.
      \end{itemize}
    \item Esto permite recorrer \textbf{todos los subconjuntos} con un for.
  \end{itemize}
\end{frame}

\begin{frame}[fragile]{Recorriendo todos los subconjuntos}
  \begin{block}{Esquema básico}
  \begin{lstlisting}[style=customc]
// mk representa un subconjunto de {0,1,2,...,n-1}
forn(mk, 1 << n) {
  forn(i, n) {
    if (mk & (1 << i)) {
      // i es elemento del subconjunto mk
    }
  }
}\end{lstlisting}
  \end{block}
  \begin{itemize}
    \item Costo: \(O(2^n \cdot n)\), aceptable para \(n \le 20\) aprox.
  \end{itemize}
\end{frame}

\begin{frame}[fragile]{Ejemplo: \href{https://cses.fi/problemset/task/1623}{Apple Division (CSES 1623)}}
  \begin{itemize}
    \item Dado un arreglo, dividirlo en dos grupos con suma lo más parecida posible.
    \item Fuerza bruta: probamos todos los subconjuntos como \textbf{un lado} de la partición.
  \end{itemize}
  \begin{block}{Implementación típica}
  \begin{lstlisting}[style=customc]
ll best = LLONG_MAX;
forn(mk, 1 << n) {
  ll s = 0;
  forn(i, n) if (mk & (1 << i)) s += a[i];
  ll resto = total - s;
  ll diff = abs(s - resto);
  best = min(best, diff);
}
cout << best << endl;\end{lstlisting}
  \end{block}
\end{frame}

\section{\texttt{next\_permutation}}

\begin{frame}[fragile]{Fuerza bruta sobre permutaciones}
  \begin{itemize}
    \item A veces queremos probar todos los órdenes posibles de un arreglo.
    \item En C++ la función \texttt{std::next\_permutation} genera la siguiente permutación lexicográfica.
  \end{itemize}
  \begin{block}{Uso típico}
\begin{lstlisting}[style=customc]
sort(v.begin(), v.end());
do {
  // usar la permutación actual de v
} while (next_permutation(v.begin(), v.end()));
\end{lstlisting}
  \end{block}
  \begin{itemize}
    \item Complejidad: \(O(n! \cdot \text{costo de procesar una permutación})\).
    \item Usable para \(n \le 10\) aprox.
  \end{itemize}
\end{frame}

\begin{frame}[fragile]{Ejemplo: Codeforces 431B}
  \begin{itemize}
    \item \href{https://codeforces.com/contest/431/problem/B}{Problema 431B - Shower Line}.
    \item 5 estudiantes se ordenan en una fila para usar la ducha.
    \item Mientras esperan, conversan por pares (1 con 2, 3 con 4, el 5 queda solo).
    \item Cada vez que alguien entra, la fila se acorta y se vuelven a formar parejas.
    \item Cada conversación entre \(i\) y \(j\) suma felicidad según una matriz \(g[i][j]\).
    \item Objetivo: encontrar el orden inicial que maximiza la felicidad total.
  \end{itemize}
\end{frame}

\begin{frame}[fragile]{Solución por fuerza bruta}
  \begin{itemize}
    \item Como \(n = 5\), alcanza con probar las \(5!\) permutaciones.
    \item Costo total \(O(n! \cdot n^2)\), más que suficiente para este límite.
  \end{itemize}
  \begin{block}{Implementación}
  \begin{lstlisting}[style=customc]
int const n = 5;
int g[n][n], p[] = {0, 1, 2, 3, 4};
int main() {
  forn(i, n) forn(j, n) cin >> g[i][j];
  ll best = 0;
  do {
    ll happiness = 0;
    forn(i, n) for (int j = i+1; j < n; j += 2)
      happiness += g[p[j-1]][p[j]] + g[p[j]][p[j-1]];
    best = max(best, happiness);
  } while (next_permutation(p, p + n));
  cout << best << endl;
}\end{lstlisting}
  \end{block}
\end{frame}

\begin{frame}{Más allá de la fuerza bruta pura}
  \begin{itemize}
    \item Para \(n\) un poco más grande podemos intentar:
      \begin{itemize}
        \item Optimizar el cálculo interno a \(O(n)\) por permutación.
        \item O usar DP sobre subconjuntos (\(O(2^n \cdot n)\)) cuando el problema lo permite.
      \end{itemize}
    \item Idea general: \textbf{aprovechar estructura} del problema para reducir el espacio de búsqueda.
  \end{itemize}
\end{frame}

\section{Fuerza bruta recursiva: N reinas}

\begin{frame}{Fuerza bruta recursiva}
  \begin{itemize}
    \item Otra forma de fuerza bruta es construir soluciones \textbf{paso a paso} recursivamente.
    \item Problema clásico: \href{https://cses.fi/problemset/task/1624}{N reinas (Chessboard and Queens, CSES 1624)}.
    \item Queremos colocar \(N\) reinas en un tablero \(N \times N\) sin que se ataquen entre sí.
  \end{itemize}
\end{frame}

\begin{frame}[fragile]{Versión muy bruta}
  \begin{block}{Recorriendo todas las configuraciones}
  \begin{lstlisting}[style=customc]
int const maxn = 10;
int N; // tamaño del tablero
bool ocupado[maxn * maxn];
bool cumple_condiciones(); // false si dos reinas se atacan entre si

bool reinas(int n) {
  if (n == 0) return cumple_condiciones();
  forn(p, N * N - (n - 1)) {
    if (ocupado[p]) continue;
    ocupado[p] = true;
    if (reinas(n - 1)) return true;
    ocupado[p] = false;
  }
  return false;
}
  \end{lstlisting}
  \end{block}
  \begin{itemize}
    \item Correcto pero recorre \textbf{muchas configuraciones equivalentes}.
  \end{itemize}
\end{frame}

\begin{frame}{Evitando duplicados innecesarios}
  \begin{itemize}
    \item La versión anterior considera diferentes órdenes para colocar reinas en las mismas casillas.
    \item Ejemplo: trata como distintos
  \end{itemize}
  \vspace{0.3cm}
  \begin{center}
  \begin{tabular}{cc}
  \begin{tabular}{|c|c|}
    \hline
    1 & \phantom{X} \\ \hline
    \phantom{X} & 2 \\ \hline
  \end{tabular}
  &
  \hspace{1cm}
  \begin{tabular}{|c|c|}
    \hline
    2 & \phantom{X} \\ \hline
    \phantom{X} & 1 \\ \hline
  \end{tabular}
  \end{tabular}
  \end{center}
  \vspace{0.3cm}
  \begin{itemize}
    \item Pero ambas configuraciones tienen una reina en la casilla 0 y otra en la casilla 3: son equivalentes.
    \item Podemos imponer un \textbf{orden} para no repetir subconjuntos.
  \end{itemize}
\end{frame}

\begin{frame}[fragile]{Reinas en orden creciente de posición}
  \begin{block}{Forzando orden en las posiciones}
  \begin{lstlisting}[style=customc]
bool reinas(int n, int last) {
  if (n == 0) return cumple_condiciones();
  forr(p, last + 1, N * N - (n - 1)) {
    ocupado[p] = true;
    if (reinas(n - 1, p)) return true;
    ocupado[p] = false;
  }
  return false;
}
  \end{lstlisting}
  \end{block}
  \begin{itemize}
    \item Ahora cada conjunto de casillas se visita una sola vez.
  \end{itemize}
\end{frame}

\begin{frame}[fragile]{Una reina por columna}
  \begin{itemize}
    \item En una solución válida solo puede haber \textbf{una reina por columna}.
    \item Podemos colocar la reina \(n\)-ésima siempre en la columna \(n-1\).
  \end{itemize}
  \begin{block}{Iterando por columnas}
  \begin{lstlisting}[style=customc]
bool reinas(int n) {
  if (n == 0) return cumple_condiciones();
  forr(j, 0, N) {
    int p = (n-1) * N + j;
    ocupado[p] = true;
    if (reinas(n - 1)) return true;
    ocupado[p] = false;
  }
  return false;
}\end{lstlisting}
  \end{block}
  \begin{itemize}
    \item Solo con estas optimizaciones pasamos de horas a milisegundos para \(N = 10\).
  \end{itemize}
\end{frame}

\section{Backtracking}

\begin{frame}{Backtracking: podar temprano}
  \begin{itemize}
    \item El backtracking mejora la fuerza bruta recursiva descartando \textbf{ramas imposibles} lo antes posible.
    \item En N reinas, cuando colocamos una reina en una fila solo avanzamos si:
      \begin{itemize}
        \item la columna no está ocupada,
        \item la diagonal principal no está ocupada,
        \item la diagonal secundaria no está ocupada.
      \end{itemize}
    \item Si alguna falla, retrocedemos (\textit{backtrack}) y probamos otra opción.
  \end{itemize}
\end{frame}

\begin{frame}[fragile]{Chequeando la condición en cada paso}
  \begin{block}{Versión ingenua de backtracking}
  \begin{lstlisting}[style=customc]
bool reinas(int n) {
  if (!cumple_condiciones()) return false;
  if (n == 0) return true;
  forr(j, 0, N) {
    ocupado[(n - 1) * N + j] = true;
    if (reinas(n - 1)) return true;
    ocupado[(n - 1) * N + j] = false;
  }
  return false;
}
  \end{lstlisting}
  \end{block}
  \begin{itemize}
    \item Mejor que chequear solo al final, pero todavía podemos optimizar.
  \end{itemize}
\end{frame}

\begin{frame}[fragile]{Marcando columnas y diagonales}
  \begin{itemize}
    \item En lugar de revisar todo el tablero en cada paso, marcamos:
      \begin{itemize}
        \item columnas ocupadas (\texttt{c[j]}),
        \item diagonal principal (\texttt{d1}),
        \item diagonal secundaria (\texttt{d2}).
      \end{itemize}
    \item Si alguna está ocupada, ni siquiera hacemos la llamada recursiva.
  \end{itemize}
  \begin{block}{Backtracking eficiente}
  \begin{lstlisting}[style=customc]
bool reinas(int n) {
  if (n == 0) return true;
  forn(j, 0, N) {
    int p = (n-1) * N + j;
    if (c[j] || d1[n + j] || d2[n - j + N]) continue;
    ocupado[p] = c[j] = d1[n + j] = d2[n - j + N] = true;
    if (reinas(n - 1)) return true;
    ocupado[p] = c[j] = d1[n + j] = d2[n - j + N] = false;
  }
  return false;
}\end{lstlisting}
  \end{block}
\end{frame}

\section{Branch \& bound}

\begin{frame}{Idea de Branch \& Bound}
  \begin{itemize}
    \item Técnica para problemas de optimización combinatoria.
    \item \textbf{Branch}: dividimos el problema en subproblemas recursivos (como backtracking).
    \item \textbf{Bound}: calculamos una \textbf{cota superior} de la mejor solución posible en cada rama.
    \item Si la cota ya es peor que la mejor solución conocida, podamo la rama sin seguir explorando.
  \end{itemize}
\end{frame}

\begin{frame}[fragile]{Ejemplo: tractores en una grilla}
  \begin{itemize}
    \item Problema de \href{https://codeforces.com/contest/143/problem/E}{Codeforces 143E}.
    \item Tenemos una grilla \(N \times M\) y queremos colocar la mayor cantidad de \textbf{tractores}.
    \item Cada tractor ocupa 5 casillas en forma de T:
  \end{itemize}
  \vspace{0.2cm}
  \begin{center}
  \begin{verbatim}
###
 #
 #
  \end{verbatim}
  \end{center}
  \begin{itemize}
    \item Se pueden rotar 90 grados.
    \item Hay que dar la cantidad máxima de tractores y una configuración que lo logre.
  \end{itemize}
\end{frame}

\begin{frame}[fragile]{Ejemplo: tractores en una grilla}
  \begin{itemize}
    \item backtracking donde vamos poniendo tractores hasta que no se pueda mas
    \item fijamos un orden
    \item en cada paso hay 5 opciones
      \begin{itemize}
          \item poner tractor hacia abajo
          \item poner tractor hacia arriba
          \item poner tractor hacia la izquierda
          \item poner tractor hacia la derecha
          \item no poner tractor
      \end{itemize}
  \end{itemize}
\end{frame}



\begin{frame}[fragile]{Backtracking básico para tractores}
  \begin{block}{Probar todas las formas de poner / no poner}
  \begin{lstlisting}[style=customc]
int cnt = 0, board[maxn][maxn]; // manejado por ok(), poner(), sacar()
int go(int p) {
  if (p == n * m) return cnt;
  int const i = p/m, j = p%m;
  int ans = cnt;
  forn(r, 5) if (ok(i, j, r)) {
    poner(i, j, r);
    ans = max(ans, go(p + 1));
    sacar(i, j, r);
  }
  return ans;
}\end{lstlisting}
  \end{block}
\end{frame}

\begin{frame}[fragile]{Guardando la mejor configuración}
  \begin{itemize}
    \item El enunciado pide también \textbf{recuperar} una configuración óptima.
    \item Guardamos el mejor tablero encontrado hasta ahora.
  \end{itemize}
  \begin{block}{Versión con memoria de la mejor solución}
  \begin{lstlisting}[style=customc]
int best = 0, best_board[maxn][maxn];
int go(int p) {

  ...
  
  if (ans > best) {
    best = ans;
    memcpy(best_board, board, sizeof(board));
  }

  return ans;
}
  \end{lstlisting}
  \end{block}
\end{frame}

\begin{frame}[fragile]{Podando con una cota superior}
  \begin{itemize}
    \item Lo anterior funciona, pero el algoritmo es muy lento.
    \item Idea: si incluso llenando todas las casillas vacías no podemos superar \texttt{best}, no vale la pena seguir.
  \end{itemize}
  \begin{block}{Branch \& bound}
  \begin{lstlisting}[style=customc]
int go(int p) {
  if (p == n * m) return cnt;

  int rem = 0;
  forr(q, p, n * m) rem += (tablero[q / m][q % m] == 0);
  int opt = cnt + rem / 5;
  if (opt <= best) return opt; // no se puede mejorar

  ...
}
  \end{lstlisting}
  \end{block}
\end{frame}

\section{Orden de búsqueda}

\begin{frame}{Importancia del orden de búsqueda}
  \begin{itemize}
    \item Mientras construimos soluciones recursivamente, el \textbf{orden} en que probamos opciones importa mucho.
    \item Un buen orden puede hacer que encontremos soluciones rápido y podamos podar más.
    \item En el problema anterior es importante que la opcion "no poner tractor" sea la ultima que probamos. Esto es asi la busqueda encuentra rapidamente soluciones con "bastantes" tractores y la idea branch\&bound puede podar mucho
    \item Ejemplo clásico: el paseo de un caballo en un tablero de ajedrez (Knight's tour).
  \end{itemize}
\end{frame}

\begin{frame}[fragile]{Heurística para Knight's tour}
  \begin{itemize}
    \item Dado un tablero \(N \times M\), queremos que un caballo visite todas las casillas exactamente una vez y vuelva al inicio.
    \item Backtracking básico es demasiado lento.
    \item Heurística: siempre movernos a la casilla con \textbf{menos vecinos libres}.
  \end{itemize}
  \begin{itemize}
    \item Para \(N = M = 6\):
      \begin{itemize}
        \item backtracking básico tarda \(\sim 20\) segundos;
        \item con la heurística tarda \(\sim 0{,}01\) segundos.
      \end{itemize}
    \item Para \(N = M = 8\), el backtracking básico es prácticamente imposible de correr.
    \item Con la heurística podemos llegar a \(N = M = 200\) en milisegundos.
  \end{itemize}
\end{frame}

\begin{frame}[fragile]{Orden aleatorio en N reinas}
  \begin{itemize}
    \item Otro ejemplo: N reinas con backtracking ya optimizado.
    \item Si probamos siempre las columnas en orden creciente, a veces caemos en casos muy malos.
    \item Si probamos las columnas en orden \textbf{aleatorio}, muchas instancias se resuelven muchísimo más rápido.
  \end{itemize}
  \begin{block}{Probar columnas en orden aleatorio}
  \begin{lstlisting}[style=customc]
bool go(int i) {
  if (i == n) return true;
  int idx[n];
  forn(j, n) idx[j] = j;
  random_shuffle(idx, idx + n);
  for (int j : idx) {
    if (c[j] || d1[i + j] || d2[i - j + n]) continue;
    c[j] = d1[i + j] = d2[i - j + n] = true;
    col[i] = j;
    if (go(i + 1)) return true;
    c[j] = d1[i + j] = d2[i - j + n] = false;
  }
  return false;
}
  \end{lstlisting}
  \end{block}
\end{frame}

\begin{frame}{Resumen}
  \begin{itemize}
    \item La \textbf{búsqueda exhaustiva} (fuerza bruta) es una herramienta poderosa para \(N\) chico.
    \item Técnicas clave:
      \begin{itemize}
        \item bitmasks para subconjuntos,
        \item \texttt{next\_permutation} para permutaciones,
        \item recursión con buenas representaciones del espacio de estados,
        \item backtracking y branch \& bound para podar muchas ramas,
        \item elegir bien el \textbf{orden de búsqueda}.
      \end{itemize}
  \end{itemize}
\end{frame}

\begin{frame}{Preguntas}
  \centering
  \Huge
  \textbf{Q\&A}
\end{frame}

\end{document}